\section{Concrete TEPL Rewrite Example}
\label{app:token-counting}
Transpose composition illustrates attribute synthesis without fixed ranks
or permutations. Imports are omitted. The explicit rule supplies the
permutation calculation; the inferred rule leaves the RHS field unspecified.

\begin{lstlisting}[language=TEPL,caption={Explicit transpose composition.}]
rule transpose_compose {
    (t2: transpose (t1: transpose X)) => (tr: transpose X)
    derive {
        tr.attrs = {
            permutation: gather(
                t1.attrs.permutation,
                t2.attrs.permutation)
        };
    }
}
\end{lstlisting}

\begin{lstlisting}[language=TEPL,caption={Transpose composition with inferred attributes.}]
rule transpose_compose {
    (t2: transpose (t1: transpose X)) => (tr: transpose X)
}
\end{lstlisting}

Both versions verify without additional guards under the evaluation's
algebraic policy. The inferred constructor is
\code{gather(t1.attrs.permutation, t2.attrs.permutation)}, matching the
explicit expression. For example, source permutations $[2,0,1]$ and
$[1,0,2]$ produce $[0,2,1]$.

\paragraph{Counting procedure.}
The evaluation's lexical scanner counts each identifier, literal, and
recognized multi-character operator as one token; other punctuation counts
individually. Whitespace and comments count zero. Counts cover the complete
rule, excluding imports and shared operation declarations. Each token belongs
to exactly one responsibility from Table~\ref{tab:xla-specification-size}.

\begin{table}[!htbp]
\caption{Token counts for the two displayed rules.}
\label{tab:transpose-token-counts}
\small
\begin{tabularx}{\linewidth}{@{}p{0.28\linewidth}Xrr@{}}\toprule
Responsibility & Counted portion & Explicit & Inferred\\\midrule
Matching pattern & Complete LHS pattern & 11 & 11\\
Legality checks & Shape constraints and \code{where} block; neither is present & 0 & 0\\
Attribute reasoning and syntax & Complete \code{derive} block & 26 & 0\\
Replacement construction & RHS pattern, including its node label & 6 & 6\\
Graph-management code & \code{rule}, rule name, outer braces, and \code{=>} & 5 & 5\\\midrule
\textbf{Total} & & \textbf{48} & \textbf{22}\\\bottomrule
\end{tabularx}
\end{table}

These counts describe the displayed variants, which retain node labels in
both versions and supply the explicit permutation without a descriptor
spread. All 26 saved tokens belong to the removed \code{derive} block;
generated attribute code is not counted as author-written specification.
The zero legality count means no explicit constraints are written, not that
operation legality is unchecked.
\FloatBarrier

\paragraph{Semantically supported, but verification unresolved.}
Static slice has declared semantics, but lies outside the whole-axis
completeness fragment because it selects coordinates within an axis.
Semantic support therefore permits verification attempts without
guaranteeing success. In this example, all RHS attributes are supplied,
yet the recorded verifier run cannot establish that the summed start
offset fits its integer type. The result is unresolved, not a proof
that the rewrite is invalid.

\begin{lstlisting}[language=TEPL,caption={Slice composition with supplied attributes.}]
rule slice_compose {
    X: [N]
    (s2: slice (s1: slice X)) => (s3: slice X)
    where {
        s1.attrs.strides == [1]
        && s2.attrs.strides == [1]
    }
    derive {
        s3.attrs = {
            starts: [s1.attrs.starts[0] + s2.attrs.starts[0]],
            sizes: s2.attrs.sizes,
            strides: [1]
        };
    }
}
\end{lstlisting}
